\documentclass[10pt,a4paper]{amsart}
\usepackage{amsmath,amssymb,amsthm,graphicx}
% Retained preamble for the Stacks Project chapter {CH} (shared project preamble).
% The chapter file pinned-original-source.tex only carries \input{{preamble}} on its first line,
% so no preamble text lives inside the chapter itself.  The Stacks Project uses one single shared
% preamble for every chapter: the snapshot stages it three times (sources/stacks-chow/preamble.tex,
% sources/stacks-etale-cohomology/preamble.tex and focused-seeds/AGP/sources/stacks-snapshot-20261011/
% preamble.tex) and those copies are identical in size.  The block below is copied mechanically,
% line for line, from sources/stacks-chow/preamble.tex lines 14-18 and 154-217: the Xy-pic setup and
% the native theorem declarations together with the mathematical macros.  Nothing was added, removed
% or redefined; the wrapper class and the amsmath/amssymb/amsthm/graphicx base packages are supplied
% by the assembler.
% For commutative diagrams we use Xy-pic
% Chapter: stacks-spaces-more-morphisms.tex -- More on Morphisms of Spaces.
\usepackage[all]{xy}

% We use 2cell for 2-commutative diagrams.
\xyoption{2cell}

% Theorem environments.
%
\theoremstyle{plain}
\newtheorem{theorem}[subsection]{Theorem}
\newtheorem{proposition}[subsection]{Proposition}
\newtheorem{lemma}[subsection]{Lemma}

\theoremstyle{definition}
\newtheorem{definition}[subsection]{Definition}
\newtheorem{example}[subsection]{Example}
\newtheorem{exercise}[subsection]{Exercise}
\newtheorem{situation}[subsection]{Situation}

\theoremstyle{remark}
\newtheorem{remark}[subsection]{Remark}
\newtheorem{remarks}[subsection]{Remarks}

\numberwithin{equation}{subsection}

% Macros
%
\def\lim{\mathop{\mathrm{lim}}\nolimits}
\def\colim{\mathop{\mathrm{colim}}\nolimits}
\def\Spec{\mathop{\mathrm{Spec}}}
\def\Hom{\mathop{\mathrm{Hom}}\nolimits}
\def\Ext{\mathop{\mathrm{Ext}}\nolimits}
\def\SheafHom{\mathop{\mathcal{H}\!\mathit{om}}\nolimits}
\def\SheafExt{\mathop{\mathcal{E}\!\mathit{xt}}\nolimits}
\def\Sch{\mathit{Sch}}
\def\Mor{\mathop{\mathrm{Mor}}\nolimits}
\def\Ob{\mathop{\mathrm{Ob}}\nolimits}
\def\Sh{\mathop{\mathit{Sh}}\nolimits}
\def\NL{\mathop{N\!L}\nolimits}
\def\CH{\mathop{\mathrm{CH}}\nolimits}
\def\proetale{{pro\text{-}\acute{e}tale}}
\def\etale{{\acute{e}tale}}
\def\QCoh{\mathit{QCoh}}
\def\Ker{\mathop{\mathrm{Ker}}}
\def\Im{\mathop{\mathrm{Im}}}
\def\Coker{\mathop{\mathrm{Coker}}}
\def\Coim{\mathop{\mathrm{Coim}}}

% Boxtimes
%
\DeclareMathSymbol{\boxtimes}{\mathbin}{AMSa}{"02}

%
% Macros for moduli stacks/spaces
%
\def\QCohstack{\mathcal{QC}\!\mathit{oh}}
\def\Cohstack{\mathcal{C}\!\mathit{oh}}
\def\Spacesstack{\mathcal{S}\!\mathit{paces}}
\def\Quotfunctor{\mathrm{Quot}}
\def\Hilbfunctor{\mathrm{Hilb}}
\def\Curvesstack{\mathcal{C}\!\mathit{urves}}
\def\Polarizedstack{\mathcal{P}\!\mathit{olarized}}
\def\Complexesstack{\mathcal{C}\!\mathit{omplexes}}
% \Pic is the operator that assigns to X its picard group, usage \Pic(X)
% \Picardstack_{X/B} denotes the Picard stack of X over B
% \Picardfunctor_{X/B} denotes the Picard functor of X over B
\def\Pic{\mathop{\mathrm{Pic}}\nolimits}
\def\Picardstack{\mathcal{P}\!\mathit{ic}}
\def\Picardfunctor{\mathrm{Pic}}
\def\Deformationcategory{\mathcal{D}\!\mathit{ef}}
\usepackage{cleveref}
\crefname{originallemma}{original lemma}{original lemmas}
\Crefname{originallemma}{Original lemma}{Original lemmas}
\makeatletter
\newlabel{lemma-localize-differentials}{{04CU}{source}{}{sourcepointer-1}{}}
\newlabel{lemma-localize-differentials@cref}{{[originallemma][0][]04CU}{[1][0][]source}{}{}{}}
\makeatother
\crefname{originallemma}{original lemma}{original lemmas}
\Crefname{originallemma}{Original lemma}{Original lemmas}
\makeatletter
\newlabel{morphisms-lemma-differentials-diagonal}{{08S2}{source}{}{sourcepointer-2}{}}
\newlabel{morphisms-lemma-differentials-diagonal@cref}{{[originallemma][0][]08S2}{[1][0][]source}{}{}{}}
\makeatother
\crefname{originallemma}{original lemma}{original lemmas}
\Crefname{originallemma}{Original lemma}{Original lemmas}
\makeatletter
\newlabel{spaces-properties-lemma-characterize-quasi-coherent}{{03M0}{source}{}{sourcepointer-3}{}}
\newlabel{spaces-properties-lemma-characterize-quasi-coherent@cref}{{[originallemma][0][]03M0}{[1][0][]source}{}{}{}}
\makeatother
\title{The sheaf of differentials of a morphism of algebraic spaces is quasi-coherent}
\author{The Stacks Project Authors}
\date{stacks-project:spaces-more-morphisms.tex:local-snapshot-20261011; GNU FDL 1.2 or later}
\begin{document}
\maketitle
\noindent\textit{Editorial note.} All delivered mathematical text is the author's own, copied from the pinned Stacks Project chapter file spaces-more-morphisms.tex; no statement or proof step was rewritten, reordered or removed.\par\medskip
\noindent\textit{Editorial note.} Preamble provenance, as required for this source: the first line of the pinned chapter file is \textbackslash{}input\{preamble\}, and the native preamble was not staged with the source, so no preamble text lives inside the chapter. The retained preamble used here is the project's single Stacks-wide preamble (the same 217-line / 8116-byte file that the already published chow and etale-cohomology cards use), copied mechanically line for line from sources/stacks-chow/preamble.tex lines 14-18 (Xy-pic setup) and 154-217 (theorem declarations and mathematical macros). Not one character was changed, no package was added, no macro was defined, and the excerpt serves only as preamble, never as a span coordinate.\par\medskip
\noindent\textit{Editorial note.} The same-chapter pointer is located at its \textbackslash{}label line inside the pinned chapter file; the two cross-chapter pointers are located in the project's own tag list staged in the snapshot (sources/stacks-chow/tags.txt), at the verbatim lines carrying their canonical tags.\par\medskip
\noindent\textit{Allowed standard prerequisites (editorial).}\begin{itemize}
\item Local presentation of a morphism of algebraic spaces by a diagram of schemes. Same chapter, Lemma labelled lemma-localize-differentials (source line 981, Stacks tag 04CU); invoked at proof lines 1023-1024. Supplies the diagram with surjective a and b and U, V schemes, which is what reduces the statement to the scheme case.
\item Quasi-coherence of the sheaf of differentials of a morphism of schemes. Morphisms, Lemma labelled morphisms-lemma-differentials-diagonal (canonical Stacks tag 08S2); invoked at proof line 1028. Supplies quasi-coherence of Omega\_\{U/V\} in the local model.
\item Criterion for quasi-coherence on an algebraic space by local presentation. Properties of Spaces, Lemma labelled spaces-properties-lemma-characterize-quasi-coherent (canonical Stacks tag 03M0); invoked at proof lines 1029-1031. Supplies the passage from quasi-coherence of the restriction to a presentation to quasi-coherence on X.
\end{itemize}\medskip
\ifdefined\hypertarget\hypertarget{sourcepointer-1}{}\fi
\noindent\textit{Source reference.} Original lemma 04CU: sources/stacks-spaces-more-morphisms/pinned-original-source.tex, source lines 981--981. Same-chapter pointer invoked at proof line 1024 for the presentation diagram; the identification of the restricted differentials and the final quasi-coherence criterion are retained in the counted span.\par\medskip
\ifdefined\hypertarget\hypertarget{sourcepointer-2}{}\fi
\noindent\textit{Source reference.} Original lemma 08S2: sources/stacks-chow/tags.txt, source lines 11115--11115. Cross-chapter standard prerequisite offered at proof line 1028 as one way to see quasi-coherence of Omega\_\{U/V\}; the reduction through the presentation is retained in the counted span.\par\medskip
\ifdefined\hypertarget\hypertarget{sourcepointer-3}{}\fi
\noindent\textit{Source reference.} Original lemma 03M0: sources/stacks-chow/tags.txt, source lines 4732--4732. Cross-chapter standard prerequisite invoked once at proof lines 1029-1031 for the criterion that quasi-coherence on an algebraic space may be checked on a presentation.\par\medskip
\begin{lemma}
\label{lemma-module-differentials-quasi-coherent}
Let $S$ be a scheme. Let $f : X \to Y$ be a morphism of algebraic spaces
over $S$. Then $\Omega_{X/Y}$ is a quasi-coherent $\mathcal{O}_X$-module.
\end{lemma}

\begin{proof}
Choose a diagram as in
Lemma \ref{lemma-localize-differentials}
with $a$ and $b$ surjective and $U$ and $V$ schemes.
Then we see that $\Omega_{X/Y}|_U = \Omega_{U/V}$ which is
quasi-coherent (for example by
Morphisms, Lemma \ref{morphisms-lemma-differentials-diagonal}).
Hence we conclude that $\Omega_{X/Y}$ is quasi-coherent by
Properties of Spaces,
Lemma \ref{spaces-properties-lemma-characterize-quasi-coherent}.
\end{proof}

\end{document}
